\chapter{Implementation Approaches}
\label{ch:implementation}

% Working draft - S. Vene, 2026-02-07

\section{Introduction}
\label{sec:implementation-intro}

Chapters~\ref{ch:trust-framework} and \ref{ch:governance} established patterns for agent integration and the trust mechanisms required for enterprise deployment. This chapter examines the practical question: \textit{how do we build these systems?}

The implementation landscape for agentic systems is fragmented. Organizations face choices across multiple dimensions:

\begin{itemize}
    \item \textbf{Orchestration paradigm:} How do agents coordinate and communicate?
    \item \textbf{State management:} In-memory vs. session-based vs. persistent
    \item \textbf{Tool integration:} How do agents interface with existing platform tooling?
    \item \textbf{Trust architecture:} Where do trust boundaries live?
\end{itemize}

This chapter provides a systematic analysis of these choices, evaluates the current tool ecosystem, and examines how agent systems integrate with established platform engineering infrastructure.

\section{Architectural Options for Agent Integration}
\label{sec:architectural-options}

\subsection{Deployment Topology}
\label{subsec:deployment-topology}

Agent systems can be deployed across three primary topologies, each with distinct trust and operational characteristics:

\textbf{Topology 1: Embedded Agents (In-Process)}

Agents run as threads or coroutines within a single process.

\begin{table}[htbp]
\centering
\caption{Embedded Agent Topology Characteristics}
\label{tab:embedded-topology}
\begin{tabular}{ll}
\toprule
\textbf{Property} & \textbf{Characteristic} \\
\midrule
Latency & Minimal (shared memory, no serialization) \\
Isolation & None (shared failure domain) \\
State management & In-memory, lost on crash \\
Scaling & Vertical only (single process limits) \\
Trust boundary & Process boundary only \\
\bottomrule
\end{tabular}
\end{table}

\textbf{Suitability:} Development, single-user tools, latency-critical applications where isolation is not required.

\textbf{Topology 2: Sidecar Agents (Per-Application)}

Agents run as separate processes alongside the primary application, communicating via IPC, RPC, or message queues.

\begin{table}[htbp]
\centering
\caption{Sidecar Agent Topology Characteristics}
\label{tab:sidecar-topology}
\begin{tabular}{ll}
\toprule
\textbf{Property} & \textbf{Characteristic} \\
\midrule
Latency & Low (local IPC, typically $<$ 10ms) \\
Isolation & Process-level (crash isolation) \\
State management & Independent per sidecar \\
Scaling & Horizontal with application replicas \\
Trust boundary & Process + network policy \\
\bottomrule
\end{tabular}
\end{table}

\textbf{Suitability:} Microservices environments, Kubernetes deployments, applications requiring crash isolation.

\textbf{Topology 3: Centralized Agent Service (Shared)}

A dedicated agent service handles requests from multiple applications.

\begin{table}[htbp]
\centering
\caption{Centralized Agent Service Characteristics}
\label{tab:centralized-topology}
\begin{tabular}{ll}
\toprule
\textbf{Property} & \textbf{Characteristic} \\
\midrule
Latency & Medium (network call, 10--100ms typical) \\
Isolation & Strong (separate service, tenant isolation possible) \\
State management & Centralized, persistent \\
Scaling & Independent of application scaling \\
Trust boundary & Network + IAM + tenant isolation \\
\bottomrule
\end{tabular}
\end{table}

\textbf{Suitability:} Enterprise deployments, multi-tenant environments, scenarios requiring centralized governance.

\subsection{Choosing Topology by Trust Requirements}
\label{subsec:topology-trust}

The Trust Equation from Chapter~\ref{ch:trust-framework} guides topology selection:

\begin{table}[htbp]
\centering
\caption{Topology Selection by Trust Requirements}
\label{tab:topology-trust}
\begin{tabular}{lllll}
\toprule
\textbf{Topology} & \textbf{Observability} & \textbf{Reversibility} & \textbf{Blast Radius} & \textbf{Suitable Autonomy} \\
\midrule
Embedded & Low & Low & Unbounded & Low (L1--L2) \\
Sidecar & Medium & Medium & Bounded per pod & Medium (L2--L3) \\
Centralized & High & High & Controlled & High (L3--L4) \\
\bottomrule
\end{tabular}
\end{table}

\textbf{Recommendation:} Enterprise deployments requiring L3+ autonomy should use centralized or sidecar topologies with explicit trust boundaries.

\section{Session-Based vs. In-Memory Orchestration}
\label{sec:session-vs-inmemory}

A fundamental architectural decision is how agent state is managed during execution. This choice has profound implications for trust, reliability, and operational characteristics.

\subsection{In-Memory Orchestration}
\label{subsec:inmemory-orchestration}

In-memory orchestration maintains agent state as runtime objects:

\begin{table}[htbp]
\centering
\caption{In-Memory Orchestration Characteristics}
\label{tab:inmemory-characteristics}
\begin{tabular}{ll}
\toprule
\textbf{Property} & \textbf{In-Memory Behavior} \\
\midrule
State persistence & Lost on process termination \\
Recovery & Requires checkpointing (often absent) \\
Latency & Minimal (no serialization) \\
Concurrency & Complex (shared mutable state) \\
Debugging & Requires runtime inspection \\
Audit trail & Must be explicitly implemented \\
\bottomrule
\end{tabular}
\end{table}

\textbf{Representative frameworks:} LangGraph, CrewAI, AutoGen

\textbf{Trust implications:}
\begin{enumerate}
    \item \textbf{Observability challenge:} Agent state exists only in memory; capturing decision provenance requires explicit instrumentation
    \item \textbf{Reversibility limitation:} Without checkpointing, rollback requires re-execution from the beginning
    \item \textbf{Blast radius:} Process crash loses all in-flight work; no isolation between concurrent tasks
\end{enumerate}

\subsection{Session-Based Orchestration}
\label{subsec:session-orchestration}

Session-based orchestration treats each agent interaction as a persistent session:

\begin{table}[htbp]
\centering
\caption{Session-Based Orchestration Characteristics}
\label{tab:session-characteristics}
\begin{tabular}{ll}
\toprule
\textbf{Property} & \textbf{Session-Based Behavior} \\
\midrule
State persistence & Survives process failures \\
Recovery & Resume from last checkpoint \\
Latency & Higher (persistence overhead, typically 50--200ms) \\
Concurrency & Simpler (isolated sessions) \\
Debugging & Session history available post-hoc \\
Audit trail & Native (session log is the audit trail) \\
\bottomrule
\end{tabular}
\end{table}

\textbf{Representative frameworks:} OpenClaw, Temporal (with LLM integration), custom implementations

\textbf{Trust implications:}
\begin{enumerate}
    \item \textbf{Observability native:} Session logs provide complete action history
    \item \textbf{Reversibility enabled:} Checkpoint/rollback semantics support undo
    \item \textbf{Blast radius contained:} Session isolation prevents cross-task interference
\end{enumerate}

\subsection{Comparative Analysis}
\label{subsec:comparative-analysis}

\textbf{When to use in-memory:}
\begin{itemize}
    \item Interactive applications where latency is critical
    \item Single-user tools without audit requirements
    \item Development and prototyping
    \item Stateless, idempotent tasks
\end{itemize}

\textbf{When to use session-based:}
\begin{itemize}
    \item Enterprise deployments with audit requirements
    \item Long-running workflows (hours to days)
    \item Tasks requiring human-in-the-loop approval
    \item Regulated environments (NIS2, AI Act compliance)
    \item Multi-tenant systems requiring isolation
\end{itemize}

\section{Tool Ecosystem Comparison}
\label{sec:tool-ecosystem}

\subsection{Framework Overview}
\label{subsec:framework-overview}

The multi-agent framework landscape has evolved rapidly:

\begin{table}[htbp]
\centering
\caption{Agent Framework Comparison}
\label{tab:framework-comparison}
\begin{tabular}{llll}
\toprule
\textbf{Framework} & \textbf{Primary Paradigm} & \textbf{Orchestration} & \textbf{Production Readiness} \\
\midrule
LangGraph & Graph-based & In-memory (with persistence options) & High \\
CrewAI & Role-based & In-memory & Medium \\
AutoGen & Conversational & In-memory & Low (research-focused) \\
OpenClaw & Session-based & Persistent sessions & Medium \\
\bottomrule
\end{tabular}
\end{table}

\subsection{LangGraph (LangChain)}
\label{subsec:langgraph}

LangGraph represents the most mature approach to controllable agent orchestration.

\textbf{Architecture:}
\begin{itemize}
    \item Agents modeled as nodes in a directed graph
    \item Edges define transitions (can be conditional)
    \item State passed through the graph explicitly
    \item Supports cycles (iterative refinement)
\end{itemize}

\textbf{Strengths:}
\begin{itemize}
    \item Explicit control flow (auditable, predictable)
    \item LangSmith integration for observability
    \item Checkpoint persistence available (PostgreSQL, Redis)
    \item Human-in-the-loop primitives
    \item Strong typing with Pydantic
\end{itemize}

\textbf{Limitations for Enterprise:}
\begin{itemize}
    \item No native multi-tenant isolation
    \item Identity management is application responsibility
    \item Checkpoint granularity is graph-level (not action-level)
    \item Observability requires LangSmith (SaaS or self-hosted)
\end{itemize}

\subsection{CrewAI}
\label{subsec:crewai}

CrewAI provides a higher-level abstraction modeling agents as team members with roles and goals.

\textbf{Strengths:}
\begin{itemize}
    \item Intuitive mental model (team collaboration)
    \item Lower learning curve than LangGraph
    \item Built-in task delegation and handoff
    \item Enterprise tier with management features
\end{itemize}

\textbf{Limitations for Enterprise:}
\begin{itemize}
    \item Less explicit control flow (harder to audit)
    \item In-memory state management
    \item Limited checkpoint/recovery options
    \item Enterprise features require commercial license
\end{itemize}

\subsection{AutoGen (Microsoft)}
\label{subsec:autogen}

AutoGen pioneered the conversational multi-agent paradigm, where agents coordinate through natural language messages.

\textbf{Strengths:}
\begin{itemize}
    \item Flexible, emergent coordination
    \item Good for research and experimentation
    \item Strong code execution capabilities
    \item Active research community
\end{itemize}

\textbf{Limitations for Enterprise:}
\begin{itemize}
    \item Designed for research, not production
    \item No built-in persistence or recovery
    \item Emergent behavior difficult to predict/audit
    \item Minimal security/isolation features
\end{itemize}

\textbf{Recommendation:} AutoGen is suitable for research and prototyping but should not be used for production platform engineering without significant hardening.

\subsection{Framework Comparison Summary}
\label{subsec:framework-summary}

\textbf{Key insight:} No single framework addresses all enterprise requirements. Organizations should expect to build additional infrastructure around these tools.

\section{Integration with Platform Engineering Tooling}
\label{sec:platform-integration}

The value of agent-enhanced platform engineering depends on how well agents integrate with existing infrastructure tools.

\subsection{Kubernetes Operators}
\label{subsec:kubernetes-operators}

Kubernetes operators extend the Kubernetes API with custom resources and controllers. Agent integration can occur at multiple levels:

\textbf{Pattern A: Agent as Operator Consumer}

Agents interact with existing operators via standard Kubernetes APIs.

Implementation considerations:
\begin{itemize}
    \item Agent needs kubeconfig/service account with appropriate RBAC
    \item Read operations (get, list, watch) are safe for exploration
    \item Write operations (create, patch, delete) require careful scoping
    \item Agents can generate YAML and submit via GitOps
\end{itemize}

\textbf{Pattern B: Agent-Aware Operators}

Custom operators designed for agent interaction might expose:
\begin{itemize}
    \item Natural language intent in CRD specs
    \item Explanation fields in status
    \item Suggested actions for common issues
    \item Structured decision points for agent approval
\end{itemize}

\subsection{Terraform and OpenTofu}
\label{subsec:terraform}

Infrastructure as Code tools present multiple integration points:

\textbf{Pattern A: Agent-Generated IaC}

Agents generate Terraform/OpenTofu configurations from natural language.

Trust considerations:
\begin{itemize}
    \item Generated code MUST go through review (human or agent reviewer)
    \item Agents should not have terraform apply privileges
    \item Dry-run (terraform plan) can be agent-executed for validation
\end{itemize}

\textbf{Pattern B: Agent-Assisted Plan Analysis}

Agents analyze terraform plan output for risks:
\begin{itemize}
    \item Destructive operations (destroy, replace)
    \item Sensitive resource types (databases, secrets)
    \item Drift from expected state
    \item Cost implications
\end{itemize}

\subsection{GitOps (ArgoCD, Flux)}
\label{subsec:gitops}

GitOps provides a natural trust boundary for agent operations: \textbf{Git is the approval gate}.

\textbf{The GitOps Agent Pattern:}
\begin{itemize}
    \item \textbf{Agent:} Proposes changes (generates manifests, creates PRs)
    \item \textbf{Git:} Tracks all changes with attribution
    \item \textbf{Human:} Approves high-risk changes
    \item \textbf{GitOps controller:} Applies approved changes
\end{itemize}

This pattern enforces separation of concerns---agent NEVER touches production directly.

\subsection{CI/CD (GitHub Actions, GitLab CI)}
\label{subsec:cicd}

Agents can participate in CI/CD pipelines in several ways:

\textbf{Pattern A: Agent-Generated Workflows}

Agents create or modify workflow definitions.

\textbf{Pattern B: Agent as CI Job}

Agents run as steps within existing pipelines for code review, security analysis, or test generation.

\textbf{Pattern C: Agent-Triggered Pipelines}

Agents can trigger pipelines based on external events (e.g., automated rollback on critical alert).

\textbf{Trust Boundaries in CI/CD:}

\begin{table}[htbp]
\centering
\caption{CI/CD Trust Boundaries for Agents}
\label{tab:cicd-trust}
\begin{tabular}{lll}
\toprule
\textbf{Operation} & \textbf{Trust Level} & \textbf{Agent Capability} \\
\midrule
Generate workflow YAML & Medium & With review \\
Trigger workflow & High & With constraints \\
Access secrets & Critical & Never \\
Push to protected branches & Critical & Never \\
Approve deployments & Critical & Human only \\
\bottomrule
\end{tabular}
\end{table}

\section{Enterprise Deployment Patterns}
\label{sec:enterprise-patterns}

These patterns emerged from iterative stress-testing: proposing architectures, identifying failure modes, refining until the patterns survived contact with known enterprise constraints.

\subsection{Pattern: Agent-in-Toolbox}
\label{subsec:agent-toolbox}

\textbf{Context:} Large enterprises typically organize around teams with isolated tenants---each team has their own infrastructure, GitOps pipelines, and access management. Behind PAM, teams maintain ``toolboxes'' containing applications for manual infrastructure management.

\textbf{Problem:} Where do you deploy the agent?

\textbf{Solution:} Deploy the agent INTO the existing toolbox. The agent inherits the team's existing blast radius---it can affect exactly what the team can already affect manually, nothing more.

\textbf{Benefits:}
\begin{itemize}
    \item Zero new attack surface---agent uses existing access paths
    \item Blast radius pre-contained---same as team's existing permissions
    \item Operational simplicity---one thing to deploy, not new infrastructure per team
    \item Incremental adoption---teams can enable/disable agent features independently
\end{itemize}

\subsection{Pattern: Read-Only Default}
\label{subsec:readonly-default}

\textbf{Problem:} LLM-based agents are non-deterministic. How do you get value from agent assistance without accepting the risk of autonomous modifications?

\textbf{Solution:} Agents operate in read-only mode by default. They can observe, query, analyze, and recommend---but not modify. Write capabilities are opt-in per team, requiring explicit enablement.

Technical enforcement, not just prompting:
\begin{itemize}
    \item \textbf{Kubernetes:} RBAC with read-only ClusterRole bound to agent's ServiceAccount
    \item \textbf{Databases:} Dedicated user with SELECT-only grants
    \item \textbf{Cloud:} IAM policies with ReadOnly managed policies
    \item \textbf{Git:} Token with read permissions only
\end{itemize}

\textbf{The Trust Progression:}
\begin{lstlisting}[caption={Trust Progression Levels}]
Level 0: Read-only, all domains
Level 1: Read-only + write to non-production
Level 2: Read-only + write to production with approval
Level 3: Autonomous write within defined boundaries
Level 4: Full autonomy (rarely appropriate)
\end{lstlisting}

Most teams should stabilize at Level 1--2. Level 3+ requires demonstrated track record and robust guardrails.

\subsection{Pattern: Federated NOC Integration}
\label{subsec:federated-noc}

\textbf{Problem:} How does NOC get cross-tenant incident visibility without violating tenant isolation?

\textbf{Anti-Pattern:} NOC Super-Agent with read access to all tenants---privileged network paths to every tenant, single compromise = total visibility loss.

\textbf{Solution: Push-Based Federation}

Team agents push incident summaries OUT to NOC. NOC receives pre-digested analysis without direct tenant access.

\textbf{Benefits:}
\begin{itemize}
    \item Tenant isolation preserved---no privileged inbound access
    \item Audit clean---data flow is unidirectional, logged at push point
    \item Scalable---adding teams doesn't increase NOC access scope
    \item Degraded gracefully---if one team agent is down, others still report
\end{itemize}

\subsection{Pattern: Event-Driven Triage}
\label{subsec:event-driven}

\textbf{Problem:} When should an agent start investigating?

\textbf{Solution:} Monitoring webhooks trigger agent triage on threshold violations.

\begin{lstlisting}[caption={Event-Driven Triage Routing}]
Prometheus/Datadog/etc.
    |
    +- Info: Log only (no agent)
    |
    +- Warning: Webhook -> Agent (background triage)
    |                        -> Queue summary for team
    |
    +- Critical: Webhook -> Agent (immediate triage)
                             +- Summary to NOC
                             +- Enriched context to PagerDuty
\end{lstlisting}

\subsection{Pattern: Vault-Integrated Credential Management}
\label{subsec:vault-integration}

\textbf{Problem:} Agents need credentials to access systems. Hard-coded or long-lived credentials create security risk.

\textbf{Solution:} Integrate with HashiCorp Vault (or OpenBao) for dynamic, short-lived, scoped credentials.

\begin{lstlisting}[caption={Vault-Integrated Credential Flow}]
Agent Triage Session:
    |
    +- 1. Authenticate to Vault (AppRole / K8s auth)
    |
    +- 2. Request scoped credential
    |      -> "Give me read-only kubectl for namespace X"
    |
    +- 3. Receive short-lived credential (TTL: 15 min)
    |
    +- 4. Perform triage operations
    |
    +- 5. Session ends
    |
    +- 6. Credential auto-expires (or explicit revoke)
\end{lstlisting}

\textbf{Benefits:}
\begin{itemize}
    \item No long-lived secrets in agent configuration
    \item Credential scope enforced by Vault policy, not just agent prompting
    \item Audit trail of all credential access
    \item Compromise limited---attacker gets short-lived, scoped credential at most
\end{itemize}

\subsection{Pattern: Git-Backed Agent Memory}
\label{subsec:git-memory}

\textbf{Problem:} Agents with persistent memory accumulate team-specific knowledge. How do you persist, version, and govern agent memory?

\textbf{Solution:} Agent memory stored in Git repositories---version controlled, auditable, recoverable.

\begin{lstlisting}[caption={Git-Backed Memory Structure}]
team-alpha-agent-memory/
  +-- incidents/
  |     +-- 2026-01-15-db-pool-exhaustion.md
  |     +-- 2026-02-03-auth-cascade.md
  +-- learnings/
  |     +-- service-quirks.md
  |     +-- deployment-patterns.md
  +-- runbooks/
  |     +-- generated-from-incidents.md
  +-- MEMORY.md (active context)
  +-- .git/
\end{lstlisting}

\textbf{Benefits:}
\begin{itemize}
    \item Memory survives agent restarts, upgrades, migrations
    \item Version history enables ``what did agent believe when incident X happened?''
    \item Human review possible for significant learnings
    \item Backup is standard git workflow
\end{itemize}

\subsection{Pattern: ChatOps Delivery}
\label{subsec:chatops}

\textbf{Problem:} Teams live in Slack/Teams. Dashboards require context switching.

\textbf{Solution:} Agent posts triage summaries and incident updates to team channels.

\textbf{Why Output-Only:}

The ``@agent restart the pod'' trap:
\begin{itemize}
    \item If agent accepts commands in chat, users will issue them
    \item Chat context is ambiguous, easy to misinterpret
    \item No approval workflow in chat interface
    \item Audit trail unclear
\end{itemize}

Instead: Agent provides analysis in chat. If human wants agent to act, they go to toolbox UI where proper guardrails exist.

\subsection{Enterprise Pattern Summary}
\label{subsec:pattern-summary}

\begin{table}[htbp]
\centering
\caption{Enterprise Pattern Language Summary}
\label{tab:enterprise-patterns}
\begin{tabular}{ll}
\toprule
\textbf{Pattern} & \textbf{Core Principle} \\
\midrule
Agent-in-Toolbox & Inherit existing access boundaries, don't create new ones \\
Read-Only Default & Capability restriction as risk management \\
Federated NOC & Push model preserves isolation \\
Event-Driven Triage & Agents activate on events, not polling \\
Cascading Failure Coordination & Hierarchy prevents thundering herd \\
Vault Credentials & Dynamic, scoped, short-lived secrets \\
Git-Backed Memory & Version control for institutional knowledge \\
ChatOps Delivery & Output where teams work, commands where guardrails exist \\
\bottomrule
\end{tabular}
\end{table}

These patterns emerged from stress-testing proposals against real constraints. They are not theoretically optimal---they are practically survivable.

\section{Summary}
\label{sec:implementation-summary}

This chapter has examined the practical dimensions of implementing agent-enhanced platform engineering:

\begin{itemize}
    \item \textbf{Topology choices} that determine trust boundaries and operational characteristics
    \item \textbf{State management approaches} with different tradeoffs for reliability and auditability
    \item \textbf{Tool ecosystem} gaps that organizations must address
    \item \textbf{Integration patterns} with established platform engineering infrastructure
    \item \textbf{Enterprise deployment patterns} that survive contact with real organizational constraints
\end{itemize}

The trust framework from Chapter~\ref{ch:trust-framework} is scale-invariant---the principles apply whether you're a startup with a single jump host or an enterprise with 100+ teams. What changes is the \textit{implementation}, not the \textit{principle}.
